\documentclass[8pt]{extarticle}

% --- Core math packages
\usepackage{amsmath, amssymb}

% --- Page + spacing tweaks to fit exactly two pages
\usepackage[margin=0.45in]{geometry}
\usepackage{microtype}
\setlength{\parskip}{3pt}
\setlength{\parindent}{0pt}
\linespread{0.98}

% --- Compact display math spacing
\makeatletter
\g@addto@macro\normalsize{%
  \setlength\abovedisplayskip{6pt}%
  \setlength\belowdisplayskip{6pt}%
  \setlength\abovedisplayshortskip{4pt}%
  \setlength\belowdisplayshortskip{4pt}%
}
\makeatother

% --- Lists tighter (keeps your content intact)
\usepackage{enumitem}
\setlist{nosep,leftmargin=*}

% --- Boxing + columns
\usepackage[most]{tcolorbox}
\usepackage{multicol}
\setlength{\columnsep}{0.35em}

\tcbset{
  colback=white,
  colframe=black!70,
  boxrule=0.5pt,
  arc=1.5pt,
  left=4pt,right=4pt,top=3pt,bottom=3pt,
  before skip=4pt, after skip=4pt,
  enhanced,
  breakable, % allow long boxes to cross a page if needed
  title style={font=\bfseries}
}

\begin{document}
\small

\begin{center}
  \Large\textbf{Math 262 -- Solutions to Practice Problems}\\[-2pt]
  \normalsize (Every original question and solution; each problem in its own box; condensed to two pages)
\end{center}

\begin{multicols}{2}
\raggedcolumns

% =========================================================
% Problem 1
% =========================================================
\begin{tcolorbox}[title=Problem 1 — Original Statement and Solution]
\noindent
\textbf{1.} Let \( T:\mathbb{R}^3 \to \mathbb{R}^2 \) be given by
\[
T\!\left(
\begin{bmatrix}x_{1}\\x_{2}\\x_{3}\end{bmatrix}
\right)
=
\begin{bmatrix}
x_{1}+x_{2}\\
x_{2}+x_{3}
\end{bmatrix}.
\]

\begin{enumerate}
\item[(a)] Show that \(T\) is a linear transformation.
\item[(b)] Find the matrix of \(T\).
\end{enumerate}
\begin{enumerate}
\item
\[
a^{(*)} \; T\!\left(\begin{bmatrix} x_{1}\\ x_{2}\\ x_{3} \end{bmatrix} +
\begin{bmatrix} y_{1}\\ y_{2}\\ y_{3} \end{bmatrix}\right)
= T\!\left(\begin{bmatrix} x_{1}+y_{1}\\ x_{2}+y_{2}\\ x_{3}+y_{3} \end{bmatrix}\right)
\]

\[
= \begin{bmatrix} x_{1}+y_{2}\\ x_{2}+y_{3} \end{bmatrix}
- \begin{bmatrix} (x_{1}+y_{1})+(x_{2}+y_{2})\\ (x_{2}+y_{2})+(x_{3}+y_{3}) \end{bmatrix}
\]

\[
= \begin{bmatrix} x_{1}+y_{2}\\ x_{2}+y_{3} \end{bmatrix}
- \begin{bmatrix} x_{1}+x_{2}+y_{1}+y_{2}\\ x_{2}+x_{3}+y_{2}+y_{3} \end{bmatrix}
\]

\[
= T\!\left(\begin{bmatrix} x_{1}\\ x_{2}\\ x_{3} \end{bmatrix}\right)
+ T\!\left(\begin{bmatrix} y_{1}\\ y_{2}\\ y_{3} \end{bmatrix}\right)
\]

\medskip
(i) For $r \in \mathbb{R}$

\[
T\!\left(r\begin{bmatrix} x_{1}\\ x_{2}\\ x_{3} \end{bmatrix}\right)
= \begin{bmatrix} r x_{1}+r x_{2}\\ r x_{2}+r x_{3} \end{bmatrix}
- r \begin{bmatrix} x_{1}+x_{2}\\ x_{2}+x_{3} \end{bmatrix}
= r \, T\!\left(\begin{bmatrix} x_{1}\\ x_{2}\\ x_{3} \end{bmatrix}\right)
\]

So $T$ is linear.

\bigskip
(ii) To find the matrix of $T$ (call it $A$):

\[
T(e_{1}) = T\!\left(\begin{bmatrix}1\\0\\0\end{bmatrix}\right)
= \begin{bmatrix}1\\0\end{bmatrix}
\]

\[
T(e_{2}) = T\!\left(\begin{bmatrix}0\\1\\0\end{bmatrix}\right)
= \begin{bmatrix}0+1\\1+0\end{bmatrix}
= \begin{bmatrix}1\\1\end{bmatrix}
\]

\[
T(e_{3}) = T\!\left(\begin{bmatrix}0\\0\\1\end{bmatrix}\right)
= \begin{bmatrix}0+0\\0+1\end{bmatrix}
= \begin{bmatrix}0\\1\end{bmatrix}
\]

\[
A = \begin{bmatrix}
1 & 1 & 0\\
0 & 1 & 1
\end{bmatrix}
\]
\noindent
\textbf{2.} Given a linear system \(A\vec{x}=\vec{b}\), suppose the reduced row echelon form (rref) of the augmented matrix is
\[
\left[
\begin{array}{cccc|c}
1 & 0 & 0 & 1 & 0\\
0 & 1 & 0 & 2 & 1\\
0 & 0 & 1 & 3 & 2\\
0 & 0 & 0 & 0 & 0
\end{array}
\right].
\]
Give the general solution to the system.
\item
\[
\begin{cases}
x_{1} + x_{4} = 0\\
x_{2} + 2x_{4} = 1\\
x_{3} + 3x_{4} = 2
\end{cases}
\quad \Rightarrow \quad
\begin{aligned}
x_{1} &= -x_{4}\\
x_{2} &= 2x_{4}+1\\
x_{3} &= -3x_{4}+2
\end{aligned}
\quad \text{letting } x_{4} = r
\]

\[
\text{The general solution: } 
\begin{bmatrix}
x_{1}\\
x_{2}\\
x_{3}\\
x_{4}
\end{bmatrix}
=
\begin{bmatrix}
-r\\
2r+1\\
-3r+2\\
r
\end{bmatrix}
\]
\end{enumerate}
\end{tcolorbox}

% =========================================================
% Problem 3
% =========================================================
\begin{tcolorbox}[title=Problem 3 — Original Statement and Solution]
\noindent
\textbf{Problem 3 (Original Statement):}  

Given the linear system
\[
\begin{cases}
x_{1} - x_{3} = 1,\\[4pt]
2x_{1} - 2x_{2} = 1,\\[4pt]
3x_{1} - 4x_{2} + x_{3} = a,
\end{cases}
\]
find (if possible) the values of \(a\) for which the system has  
\begin{enumerate}
\item[(a)] no solution,  
\item[(b)] exactly one solution,  
\item[(c)] infinitely many solutions.  
\end{enumerate}
Also, in the case when there are infinitely many solutions, give the general solution.

\bigskip
\noindent
\textbf{Solution:}

\[
\begin{bmatrix}
1 & 0 & -1 & | & 1\\
-2 & 2 & 0 & | & 1\\
3 & -4 & 1 & | & a
\end{bmatrix}
\;\to\;
\begin{bmatrix}
1 & 0 & -1 & | & 1\\
0 & -2 & 2 & | & -1\\
0 & -4 & 4 & | & a-3
\end{bmatrix}
\;\to\;
\begin{bmatrix}
1 & 0 & -1 & | & 1\\
0 & 1 & -1 & | & \tfrac{1}{2}\\
0 & 0 & 0 & | & a-1
\end{bmatrix}
\]

\[
\text{(a) No solution if } a \ne 1
\]

\[
\text{(b) There is never exactly one solution.}
\]

\[
\text{(c) Infinitely many solutions when } a = 1
\]

In that case,
\[
x_{1} = 1 + r, \quad
x_{2} = \tfrac{1}{2} + r, \quad
x_{3} = r \quad (\text{letting } x_{3} = r).
\]

\[
\begin{bmatrix}
x_{1}\\
x_{2}\\
x_{3}
\end{bmatrix}
=
\begin{bmatrix}
1 + r\\
\tfrac{1}{2} + r\\
r
\end{bmatrix}
\]
\end{tcolorbox}

% =========================================================
% Problem 4
% =========================================================
\begin{tcolorbox}[title=Problem 4 — Original Statement and Solution]
\noindent
In \(\mathbb{R}^{2}\), let
\[
\vec{u} =
\begin{bmatrix}
\dfrac{1}{\sqrt{10}}\\[6pt]
\dfrac{3}{\sqrt{10}}
\end{bmatrix}.
\]
\begin{enumerate}
\item[(a)] Show that \(\vec{u}\) is a unit vector.  
\item[(b)] Let \(L\) be the line spanned by \(\vec{u}\), and let
\[
\vec{x} =
\begin{bmatrix}
-1\\
4
\end{bmatrix}.
\]
Find the reflection of \(\vec{x}\) through \(L\), that is,
\[
\mathrm{ref}_{L}(\vec{x}).
\]
\end{enumerate}

\bigskip
\noindent
\textbf{Solution:}

\[
|\vec{u}| = \sqrt{\left(\frac{1}{\sqrt{10}}\right)^{2} + \left(\frac{3}{\sqrt{10}}\right)^{2}}
= \sqrt{\frac{1}{10} + \frac{9}{10}}
= \sqrt{1} = 1
\]

\[
\mathrm{proj}_{L}(\vec{x})
= \frac{\vec{x} \cdot \vec{u}}{\vec{u} \cdot \vec{u}} \, \vec{u}
\]

\[
\vec{x} \cdot \vec{u}
= \begin{bmatrix}-1\\4\end{bmatrix} \cdot
\begin{bmatrix}\tfrac{1}{\sqrt{10}}\\[4pt]\tfrac{3}{\sqrt{10}}\end{bmatrix}
= -\tfrac{1}{\sqrt{10}} + \tfrac{12}{\sqrt{10}}
= \tfrac{11}{\sqrt{10}}
\]

\[
\vec{u} \cdot \vec{u} = 1 \quad (\text{since } \vec{u} \text{ is a unit vector})
\]

\[
\mathrm{proj}_{L}(\vec{x})
= \tfrac{11}{\sqrt{10}}
\begin{bmatrix}
\tfrac{1}{\sqrt{10}}\\[4pt]
\tfrac{3}{\sqrt{10}}
\end{bmatrix}
= \begin{bmatrix}\tfrac{11}{10}\\[4pt]\tfrac{33}{10}\end{bmatrix}
\]

\[
\mathrm{ref}_{L}(\vec{x}) = 2\,\mathrm{proj}_{L}(\vec{x}) - \vec{x}
\]
\end{tcolorbox}

% =========================================================
% Problem 5
% =========================================================
\begin{tcolorbox}[title=Problem 5 — Original Statement and Solution]
\noindent
\textbf{Original Problem:}  
In \(\mathbb{R}^{3}\), let \(L\) be the line spanned by the vector
\[
\vec{w} =
\begin{bmatrix}
\dfrac{1}{\sqrt{2}}\\[6pt]
0\\[6pt]
\dfrac{1}{\sqrt{2}}
\end{bmatrix},
\qquad
\vec{x} =
\begin{bmatrix}
1\\
1\\
1
\end{bmatrix}.
\]
Verify that
\[
\vec{x} - \mathrm{proj}_{L}(\vec{x})
\]
is orthogonal to \(\vec{w}\), that is, their dot product is zero.

\bigskip
\noindent
\textbf{Solution:}

\[
2\begin{bmatrix}
\frac{11}{10}\\[4pt]
\frac{33}{10}
\end{bmatrix}
-
\begin{bmatrix}
-1\\[4pt]
4
\end{bmatrix}
=
\begin{bmatrix}
\frac{22}{10}\\[4pt]
\frac{66}{10}
\end{bmatrix}
-
\begin{bmatrix}
-1\\[4pt]
4
\end{bmatrix}
=
\begin{bmatrix}
\frac{22}{10}+1\\[4pt]
\frac{66}{10}-4
\end{bmatrix}
=
\begin{bmatrix}
\frac{16}{5}\\[4pt]
\frac{13}{5}
\end{bmatrix}
\]

\medskip
\noindent
\textit{Note: You could also use the matrix formula for this one.}

\bigskip
\[
\mathrm{proj}_{L}(\vec{x})
= \frac{\vec{x}\cdot \vec{w}}{\vec{w}\cdot \vec{w}}\,\vec{w}
= \frac{2}{\sqrt{2}}
\begin{bmatrix}
\frac{1}{\sqrt{2}}\\[4pt]
\frac{1}{\sqrt{2}}
\end{bmatrix}
= \begin{bmatrix}1\\1\end{bmatrix}.
\]

\[
\vec{x} - \mathrm{proj}_{L}(\vec{x})
= \begin{bmatrix}1\\1\\1\end{bmatrix}
- \begin{bmatrix}0\\1\\0\end{bmatrix}
= \begin{bmatrix}1\\0\\1\end{bmatrix}.
\]

\[
\text{Finally, }
\begin{bmatrix}1\\0\\1\end{bmatrix}
\cdot
\begin{bmatrix}\tfrac{1}{\sqrt{2}}\\[4pt]0\\[4pt]\tfrac{1}{\sqrt{2}}\end{bmatrix}
= 0,
\quad\text{so } \vec{x} - \mathrm{proj}_{L}(\vec{x}) \text{ is orthogonal to } \vec{w} \text{ as claimed.}
\]
\end{tcolorbox}

% =========================================================
% Problem 6
% =========================================================
\begin{tcolorbox}[title=Problem 6 — Original Statement and Solution]
\noindent
\textbf{Problem 6 (Original Statement):}  

Define \( T:\mathbb{R}^{2} \to \mathbb{R}^{2} \) by
\[
T\!\left(
\begin{bmatrix} x_{1}\\ x_{2} \end{bmatrix}
\right)
=
\begin{bmatrix}
x_{1} - x_{2}\\
x_{1} + x_{2}
\end{bmatrix}.
\]

You may assume \(T\) is a linear transformation.  
\begin{enumerate}
\item[(a)] Find \( A \), the matrix representation of \(T\).
\item[(b)] Use \(A\) to explain why \(T\) is invertible.
\item[(c)] Use \(A\) to find the rule for \(T^{-1}\).
\end{enumerate}

\bigskip
\noindent
\textbf{Solution:}

\[
T(e_{1}) = \begin{bmatrix}1\\1\end{bmatrix},
\qquad
T(e_{2}) = \begin{bmatrix}-1\\1\end{bmatrix},
\qquad
A = \begin{bmatrix}1 & -1\\ 1 & 1\end{bmatrix}.
\]

\[
\det A = (1)(1) - (-1)(1) = 2 \quad \Rightarrow \quad A \text{ is invertible.}
\]
Since \(T = T_{A}\), it follows that \(T\) is invertible and
\[
T^{-1} = T_{A^{-1}}.
\]

\[
A^{-1} = \frac{1}{2}
\begin{bmatrix}
1 & 1\\
-1 & 1
\end{bmatrix}
=
\begin{bmatrix}
\tfrac{1}{2} & \tfrac{1}{2}\\
-\tfrac{1}{2} & \tfrac{1}{2}
\end{bmatrix}.
\]

Thus
\[
T^{-1}
\begin{bmatrix}
x_{1}\\ x_{2}
\end{bmatrix}
=
A^{-1}
\begin{bmatrix}
x_{1}\\ x_{2}
\end{bmatrix}
=
\begin{bmatrix}
\tfrac{1}{2}(x_{1}+x_{2})\\[6pt]
\tfrac{1}{2}(-x_{1}+x_{2})
\end{bmatrix}.
\]
\end{tcolorbox}

% =========================================================
% Problem 7
% =========================================================
\begin{tcolorbox}[title=Problem 7 — Original Statement and Solution]
\noindent
Let \(\theta\) be any angle and let \(A\) be the matrix which represents the counterclockwise rotation through \(\theta\).  
Prove that for any vector
\[
\vec{u} =
\begin{bmatrix} x\\ y \end{bmatrix} \in \mathbb{R}^{2},
\quad
|A\vec{u}| = |\vec{u}|
\]
—that is, rotations preserve length.

\bigskip
\noindent
\textbf{Solution:}

\[
A =
\begin{bmatrix}
\cos\theta & -\sin\theta\\
\sin\theta & \cos\theta
\end{bmatrix},
\qquad
\vec{u} =
\begin{bmatrix} x\\ y \end{bmatrix},
\qquad
A\vec{u} =
\begin{bmatrix}
x\cos\theta - y\sin\theta\\
x\sin\theta + y\cos\theta
\end{bmatrix}.
\]

\[
|A\vec{u}|
= \sqrt{(x\cos\theta - y\sin\theta)^{2} + (x\sin\theta + y\cos\theta)^{2}}
\]

Expand:
\[
= \sqrt{
x^{2}\cos^{2}\theta - 2xy\cos\theta\sin\theta + y^{2}\sin^{2}\theta
+ x^{2}\sin^{2}\theta + 2xy\sin\theta\cos\theta + y^{2}\cos^{2}\theta
}
\]

\[
= \sqrt{
x^{2}(\cos^{2}\theta+\sin^{2}\theta) + y^{2}(\sin^{2}\theta+\cos^{2}\theta)
}
= \sqrt{x^{2} + y^{2}}
= |\vec{u}|.
\]

Hence rotations preserve length.
\end{tcolorbox}

% =========================================================
% Problem 8
% =========================================================
\begin{tcolorbox}[title=Problem 8 — Original Statement and Solution]
\noindent
Fill in the blanks for an invertible \( n \times n \) matrix \( A \):
\begin{enumerate}
\item[(a)] \(\mathrm{rank}(A) = \underline{\hspace{1cm}}\)
\item[(b)] The homogeneous system \( A\vec{x} = \vec{0} \) has only the \underline{\hspace{2cm}} solution.
\item[(c)] \( A \) is row-reducible to the \underline{\hspace{2cm}} matrix.
\end{enumerate}

\bigskip
\noindent
\textbf{Solution:}
\[
\mathrm{rank}(A) = n, \qquad
\text{homogeneous system has only the trivial solution,} \qquad
A \text{ is row-reducible to the identity matrix.}
\]
\end{tcolorbox}

% =========================================================
% Problem 9
% =========================================================
\begin{tcolorbox}[title=Problem 9 — Original Statement and Solution]
\noindent
Let
\[
A =
\begin{bmatrix}
1 & 1 & 0 & 1 & 0\\
2 & 3 & 4 & 0 & 1\\
3 & 4 & 0 & 1 & 5\\
6 & 8 & 4 & 2 & 6
\end{bmatrix}.
\]
Give the reduced row echelon form of \(A\).

\bigskip
\noindent
\textbf{Solution:}

\[
\begin{bmatrix}
1 & 1 & 0 & 1 & 0\\
2 & 3 & 4 & 0 & 1\\
3 & 4 & 0 & 1 & 5\\
6 & 8 & 4 & 2 & 6
\end{bmatrix}
\;\to\;
\begin{bmatrix}
1 & 0 & -4 & -1\\
0 & 1 & 0 & 5\\
0 & 0 & 4 & 4\\
0 & 0 & 4 & 4
\end{bmatrix}
\]

\[
\to\;
\begin{bmatrix}
1 & 0 & -4 & -1\\
0 & 1 & 0 & 5\\
0 & 0 & 1 & -1\\
0 & 0 & 0 & 0
\end{bmatrix}.
\]

So the reduced row echelon form of \(A\) is
\[
\boxed{
\begin{bmatrix}
1 & 0 & 0 & -5\\
0 & 1 & 0 & 5\\
0 & 0 & 1 & -1\\
0 & 0 & 0 & 0
\end{bmatrix}}
\]
\end{tcolorbox}

% =========================================================
% Problem 10(a)
% =========================================================
\begin{tcolorbox}[title=Problem 10(a) — Original Statement and Solution]
\noindent
(a) Let
\[
A = \begin{bmatrix}1 & 2\\ 2 & 1\end{bmatrix}.
\]
Prove that if
\[
AB = \begin{bmatrix}0 & 0\\ 0 & 0\end{bmatrix},
\]
then
\[
B = \begin{bmatrix}0 & 0\\ 0 & 0\end{bmatrix}.
\]

\bigskip
\noindent
\textbf{Solution:}

\[
A = \begin{bmatrix}1 & 2\\ 2 & 1\end{bmatrix},
\quad
B = \begin{bmatrix}a & b\\ c & d\end{bmatrix},
\quad
AB =
\begin{bmatrix}
a + 2c & b + 2d\\
2a + c & 2b + d
\end{bmatrix}.
\]

If \(AB = \begin{bmatrix}0 & 0\\ 0 & 0\end{bmatrix}\), then
\[
a + 2c = 0, \quad 2a + c = 0
\quad\Rightarrow\quad a = c = 0,
\]
and
\[
b + 2d = 0, \quad 2b + d = 0
\quad\Rightarrow\quad b = d = 0.
\]

Hence
\[
B = \begin{bmatrix}0 & 0\\ 0 & 0\end{bmatrix}.
\]
\end{tcolorbox}

% =========================================================
% Problem 10(b)
% =========================================================
\begin{tcolorbox}[title=Problem 10(b) — Original Statement and Solution]
\noindent
Let
\[
C =
\begin{bmatrix}
0 & 1\\
0 & 1
\end{bmatrix}.
\]
Give the general form of a matrix \(D\) such that
\[
DC =
\begin{bmatrix}
0 & 0\\
0 & 0
\end{bmatrix}.
\]

\bigskip
\noindent
\textbf{Solution:}

Let
\[
D =
\begin{bmatrix}
a & b\\
c & d
\end{bmatrix}.
\]
Then
\[
DC =
\begin{bmatrix}
c & d\\
c & d
\end{bmatrix}.
\]

For \(DC = \begin{bmatrix}0 & 0\\ 0 & 0\end{bmatrix}\), we must have
\[
c = 0, \qquad d = 0.
\]

Hence
\[
D =
\begin{bmatrix}
a & b\\
0 & 0
\end{bmatrix}.
\]
\end{tcolorbox}

% =========================================================
% Problem 11
% =========================================================
\begin{tcolorbox}[title=Problem 11 — Original Statement and Solution]
\noindent
In \(\mathbb{R}^{2}\), let
\[
\vec{u} = 
\begin{bmatrix}
1\\ 2
\end{bmatrix},
\qquad
\vec{v} =
\begin{bmatrix}
2\\ 1
\end{bmatrix}.
\]
Show that any vector
\[
\vec{w} =
\begin{bmatrix}
a\\ b
\end{bmatrix}
\]
can be written as a linear combination
\[
\vec{w} = r\vec{u} + s\vec{v}.
\]

\bigskip
\noindent
\textbf{Solution:}

We need to find scalars \(r\) and \(s\) such that
\[
r\begin{bmatrix}1\\2\end{bmatrix}
+ s\begin{bmatrix}2\\1\end{bmatrix}
= \begin{bmatrix}a\\b\end{bmatrix}.
\]

This gives the system
\[
\begin{cases}
r + 2s = a,\\
2r + s = b.
\end{cases}
\]

Multiply the first equation by 2:
\[
2r + 4s = 2a
\]

Subtract the second equation:
\[
(2r + 4s) - (2r + s) = 2a - b
\quad\Rightarrow\quad
3s = 2a - b
\quad\Rightarrow\quad
s = \frac{2a - b}{3}.
\]

Substitute into \(r + 2s = a\):
\[
r + 2\left(\frac{2a - b}{3}\right) = a
\quad\Rightarrow\quad
3r + 4a - 2b = 3a
\quad\Rightarrow\quad
3r = 3a - 4a + 2b = -a + 2b
\quad\Rightarrow\quad
r = \frac{-a + 2b}{3}.
\]

Hence
\[
\begin{bmatrix}a\\b\end{bmatrix}
= \left(\frac{-a + 2b}{3}\right)
\begin{bmatrix}1\\2\end{bmatrix}
+ \left(\frac{2a - b}{3}\right)
\begin{bmatrix}2\\1\end{bmatrix}.
\]
\end{tcolorbox}

% =========================================================
% Problem 12
% =========================================================
\begin{tcolorbox}[title=Problem 12 — Original Statement and Solution]
\noindent
\begin{enumerate}
\item[(a)] Is the matrix
\[
\begin{bmatrix}
1 & 2 & 3\\
4 & 5 & 6\\
7 & 8 & 9
\end{bmatrix}
\]
invertible? Why or why not?

\item[(b)] The matrix
\[
A =
\begin{bmatrix}
1 & 1 & 1\\
2 & 1 & 0\\
0 & 1 & 1
\end{bmatrix}
\]
is invertible. Find \(A^{-1}\) and use it to solve the equation
\[
A\vec{x} =
\begin{bmatrix}
1\\ -3\\ 2
\end{bmatrix}.
\]
\end{enumerate}

\bigskip
\noindent
\textbf{Solution:}

% ---------- Part (a)
\bigskip
\noindent
(a)\,
Augment the matrix with the identity:
\[
\left[
\begin{array}{ccc|ccc}
1 & 2 & 3 & 1 & 0 & 0\\
4 & 5 & 6 & 0 & 1 & 0\\
7 & 8 & 9 & 0 & 0 & 1
\end{array}
\right]
\to
\left[
\begin{array}{ccc|ccc}
1 & 2 & 3 & 1 & 0 & 0\\
0 & -3 & -6 & -4 & 1 & 0\\
0 & -6 & -12 & -3 & 0 & 1
\end{array}
\right].
\]

Further row reduction leads to a row of zeros in the left block:
\[
\left[
\begin{array}{ccc}
1 & 2 & 3\\
0 & -3 & -6\\
0 & 0 & 0
\end{array}
\right],
\]
so the matrix cannot be reduced to the identity matrix.

Therefore, the matrix is \textbf{not invertible}.

% ---------- Part (b)
\bigskip
\noindent
(b)\,
Augment \(A\) with the identity:
\[
\left[
\begin{array}{ccc|ccc}
1 & 1 & 1 & 1 & 0 & 0\\
2 & 1 & 0 & 0 & 1 & 0\\
0 & 1 & 1 & 0 & 0 & 1
\end{array}
\right]
\to
\left[
\begin{array}{ccc|ccc}
1 & 1 & 1 & 1 & 0 & 0\\
0 & -1 & -2 & -2 & 1 & 0\\
0 & 1 & 1 & 0 & 0 & 1
\end{array}
\right]
\to
\left[
\begin{array}{ccc|ccc}
1 & 0 & -1 & 1 & 0 & 1\\
0 & 1 & 2  & 2 & -1 & 0\\
0 & 0 & 1  & 2 & -1 & -1
\end{array}
\right].
\]

Hence
\[
A^{-1} =
\begin{bmatrix}
1 & 0 & -1\\
-2 & 1 & 2\\
2 & -1 & -1
\end{bmatrix}.
\]

Solve
\[
A\vec{x} =
\begin{bmatrix}
1\\ -3\\ 2
\end{bmatrix}
\quad\Rightarrow\quad
\vec{x} =
A^{-1}
\begin{bmatrix}
1\\ -3\\ 2
\end{bmatrix}
=
\begin{bmatrix}
1 & 0 & -1\\
-2 & 1 & 2\\
2 & -1 & -1
\end{bmatrix}
\begin{bmatrix}
1\\ -3\\ 2
\end{bmatrix}
=
\begin{bmatrix}
-1\\ -1\\ 3
\end{bmatrix}.
\]

Thus,
\[
\boxed{
\vec{x} =
\begin{bmatrix}
-1\\
-1\\
3
\end{bmatrix}
}.
\]
\end{tcolorbox}

% =========================================================
% Problem 13
% =========================================================
\begin{tcolorbox}[title=Problem 13 — Original Statement and Solution]
\noindent
\textbf{Problem 13 (Original Statement):}  

Let
\[
A =
\begin{bmatrix}
1 & r\\
2r & 1
\end{bmatrix}.
\]

\begin{enumerate}
\item[(a)] Find \(\det A\).
\item[(b)] For what value(s) of \(r\) is \(A\) invertible?
\item[(c)] In the cases when \(A\) is invertible, give \(A^{-1}\).
\end{enumerate}

\bigskip
\noindent
\textbf{Solution:}

\[
\det A = (1)(1) - (2r)(r) = 1 - 2r^{2}.
\]

\[
\det A = 0 \quad \Rightarrow \quad 1 - 2r^{2} = 0
\quad\Rightarrow\quad 2r^{2} = 1
\quad\Rightarrow\quad r = \pm \frac{1}{\sqrt{2}}.
\]

Thus, \(A\) is invertible for all \(r \neq \pm \frac{1}{\sqrt{2}}\).

\bigskip
\[
A^{-1}
= \frac{1}{1 - 2r^{2}}
\begin{bmatrix}
1 & -r\\
-2r & 1
\end{bmatrix}.
\]
\end{tcolorbox}

\end{multicols}
\end{document}

